% GENERATED FILE. Edit canonical lessons, then run npm run generate:books.
% canonical-source-hash: fnv1a64:7416fd4a5b12314e
% canonical-lessons: UR-C520-pahalvan, UR-C520-mugdar, UR-C520-dao, UR-C520-bhirana, UR-C520-langot

\chapter{In the Wrestling Pit}
\label{ch:ur-in-the-wrestling-pit}
\hlchaptermodality{520}

\begin{chapteropening}
\textbf{By the end of this chapter:} \emph{I can talk about wrestlers and how they train.}

Complete the last lesson of chapter 520: I can talk about wrestlers and how they train.
\end{chapteropening}

\section[\texorpdfstring{\ur{پہلوان} (pahalvān)}{پہلوان (pahalvān)}]{\ur{پہلوان} (pahalvān) — a wrestler}

\label{lesson:UR-C520-pahalvan}



\begin{quote}
Before the new one: say the Urdu for a gathering (for music or poetry), then the Urdu for to make (someone) dance.
\end{quote}

\subsection*{You'll want to know: \ur{پہلوان}}
\textbf{\ur{پہلوان}} — \emph{pahalvān} — \textquotedblleft{}a wrestler\textquotedblright{}.

\begin{quote}
\ur{پہلوان} \ur{روز} \ur{صبح} \ur{دودھ} \ur{پیتا} \ur{ہے۔}
\end{quote}

\begin{quote}
\emph{pahalvān roz subah dūdh pītā hai.}
\end{quote}

\begin{quote}
The wrestler drinks milk every morning.
\end{quote}

Masculine.

\begin{cousinweb}
\textbf{\ur{پہلوان}} (\emph{pahalvān}) is usually traced to Persian \emph{pahlavān}, \textquotedblleft{}a hero, a champion\textquotedblright{}.
\end{cousinweb}

\begin{grammarlens}[title={What you've built}]
One more word for describing what you lived and felt.
\end{grammarlens}

\subsection*{Your turn}
\begin{itemize}
\raggedright
  \item \emph{Say it:} \emph{pahalvān}
  \item \emph{Say it:} \emph{pahalvān}, once more
  \item \emph{Say it:} \emph{pahalvān roz subah dūdh pītā hai.} — The wrestler drinks milk every morning.
  \item \emph{Recall:} say the Urdu for a gathering (for music or poetry), then the Urdu for to make (someone) dance, then say \emph{pahalvān} again
\end{itemize}

\begin{tcolorbox}[breakable,colback=teal!4,colframe=teal!35!black,title={Before you move on}]
What does \ur{پہلوان} mean? (\textquotedblleft{}A wrestler\textquotedblright{}.) Say it once more.
\end{tcolorbox}
\section[\texorpdfstring{\ur{مگدر} (mugdar)}{مگدر (mugdar)}]{\ur{مگدر} (mugdar) — an Indian club (for exercise)}

\label{lesson:UR-C520-mugdar}



\begin{quote}
Before the new one: say the Urdu for to make (someone) dance, then the Urdu for a wrestler.
\end{quote}

\subsection*{You'll want to know: \ur{مگدر}}
\textbf{\ur{مگدر}} — \emph{mugdar} — \textquotedblleft{}an Indian club (for exercise)\textquotedblright{}.

\begin{quote}
\ur{پہلوان} \ur{کے} \ur{پاس} \ur{دو} \ur{بھاری} \ur{مگدر} \ur{ہیں۔}
\end{quote}

\begin{quote}
\emph{pahalvān ke pās do bhārī mugdar haiṅ.}
\end{quote}

\begin{quote}
The wrestler has two heavy clubs.
\end{quote}

Masculine.

\begin{cousinweb}
\textbf{\ur{مگدر}} (\emph{mugdar}) is inherited from Sanskrit \emph{mudgara}, \textquotedblleft{}a mallet, a club\textquotedblright{}, under which CDIAL lists Hindi \emph{mūgrā}, \textquotedblleft{}a mallet\textquotedblright{} (CDIAL 10199).
\end{cousinweb}

\begin{grammarlens}[title={What you've built}]
One more word for describing what you lived and felt.
\end{grammarlens}

\subsection*{Your turn}
\begin{itemize}
\raggedright
  \item \emph{Say it:} \emph{mugdar}
  \item \emph{Say it:} \emph{mugdar}, once more
  \item \emph{Say it:} \emph{pahalvān ke pās do bhārī mugdar haiṅ.} — The wrestler has two heavy clubs.
  \item \emph{Recall:} say the Urdu for to make (someone) dance, then the Urdu for a wrestler, then say \emph{mugdar} again
\end{itemize}

\begin{tcolorbox}[breakable,colback=teal!4,colframe=teal!35!black,title={Before you move on}]
What does \ur{مگدر} mean? (\textquotedblleft{}An Indian club (for exercise)\textquotedblright{}.) Say it once more.
\end{tcolorbox}
\section[\texorpdfstring{\ur{داؤ} (dāo)}{داؤ (dāo)}]{\ur{داؤ} (dāo) — a move (in wrestling)}

\label{lesson:UR-C520-dao}



\begin{quote}
Before the new one: say the Urdu for a wrestler, then the Urdu for an Indian club (for exercise).
\end{quote}

\subsection*{You'll want to know: \ur{داؤ}}
\textbf{\ur{داؤ}} — \emph{dāo} — \textquotedblleft{}a move (in wrestling)\textquotedblright{}.

\begin{quote}
\ur{استاد} \ur{نے} \ur{اسے} \ur{ایک} \ur{نیا} \ur{داؤ} \ur{سکھایا۔}
\end{quote}

\begin{quote}
\emph{ustād ne use ek nayā dāo sikhāyā.}
\end{quote}

\begin{quote}
His trainer taught him a new move.
\end{quote}

Masculine.

\begin{cousinweb}
\textbf{\ur{داؤ}} (\emph{dāo}) is inherited: CDIAL lists the Hindi cousin \emph{dā̃w}, \textquotedblleft{}a turn\textquotedblright{}, under Sanskrit \emph{dātu}, \textquotedblleft{}a share\textquotedblright{} (CDIAL 6258); the same entry also mentions the sense \textquotedblleft{}a move in a game\textquotedblright{}, so the step from \textquotedblleft{}a turn\textquotedblright{} to \textquotedblleft{}a move\textquotedblright{} is an old one.
\end{cousinweb}

\begin{grammarlens}[title={What you've built}]
One more word for describing what you lived and felt.
\end{grammarlens}

\subsection*{Your turn}
\begin{itemize}
\raggedright
  \item \emph{Say it:} \emph{dāo}
  \item \emph{Say it:} \emph{dāo}, once more
  \item \emph{Say it:} \emph{ustād ne use ek nayā dāo sikhāyā.} — His trainer taught him a new move.
  \item \emph{Recall:} say the Urdu for a wrestler, then the Urdu for an Indian club (for exercise), then say \emph{dāo} again
\end{itemize}

\begin{tcolorbox}[breakable,colback=teal!4,colframe=teal!35!black,title={Before you move on}]
What does \ur{داؤ} mean? (\textquotedblleft{}A move (in wrestling)\textquotedblright{}.) Say it once more.
\end{tcolorbox}
\section[\texorpdfstring{\ur{بھڑانا} (bhiṛānā)}{بھڑانا (bhiṛānā)}]{\ur{بھڑانا} (bhiṛānā) — to pit against each other}

\label{lesson:UR-C520-bhirana}



\begin{quote}
Before the new one: say the Urdu for an Indian club (for exercise), then the Urdu for a move (in wrestling).
\end{quote}

\subsection*{You'll want to know: \ur{بھڑانا}}
\textbf{\ur{بھڑانا}} — \emph{bhiṛānā} — \textquotedblleft{}to pit against each other\textquotedblright{}.

\begin{quote}
\ur{استاد} \ur{نے} \ur{دونوں} \ur{پہلوانوں} \ur{کو} \ur{بھڑایا۔}
\end{quote}

\begin{quote}
\emph{ustād ne donoṅ pahalvānoṅ ko bhiṛāyā.}
\end{quote}

\begin{quote}
The trainer pitted the two wrestlers against each other.
\end{quote}

Past, with \ur{نے} (he or she …): \ur{اس} \ur{نے} \ur{دو} \ur{لڑکوں} \ur{کو} \ur{بھڑایا}, \emph{us ne do laṛkoṅ ko bhiṛāyā}.

\begin{cousinweb}
\textbf{\ur{بھڑانا}} (\emph{bhiṛānā}) is the causative of \emph{bhiṛnā}, \textquotedblleft{}to clash\textquotedblright{}: CDIAL lists the Hindi cousins \emph{bhiṛnā} and \emph{bhiṛānā}, \textquotedblleft{}to bring together\textquotedblright{}, under \emph{*bhiṭ}, \textquotedblleft{}meet, throng\textquotedblright{}, a form it reconstructs (CDIAL 9490).
\end{cousinweb}

\begin{grammarlens}[title={What you've built}]
One more word for describing what you lived and felt.
\end{grammarlens}

\subsection*{Your turn}
\begin{itemize}
\raggedright
  \item \emph{Say it:} \emph{bhiṛānā}
  \item \emph{Say it:} \emph{bhiṛānā}, once more
  \item \emph{Say it:} \emph{ustād ne donoṅ pahalvānoṅ ko bhiṛāyā.} — The trainer pitted the two wrestlers against each other.
  \item \emph{Recall:} say the Urdu for an Indian club (for exercise), then the Urdu for a move (in wrestling), then say \emph{bhiṛānā} again
\end{itemize}

\begin{tcolorbox}[breakable,colback=teal!4,colframe=teal!35!black,title={Before you move on}]
What does \ur{بھڑانا} mean? (\textquotedblleft{}To pit against each other\textquotedblright{}.) Say it once more.
\end{tcolorbox}
\section[\texorpdfstring{\ur{لنگوٹ} (langoṭ)}{لنگوٹ (langoṭ)}]{\ur{لنگوٹ} (langoṭ) — a loincloth (a wrestler's)}

\label{lesson:UR-C520-langot}



\begin{quote}
Before the new one: say the Urdu for a move (in wrestling), then the Urdu for to pit against each other.
\end{quote}

\subsection*{You'll want to know: \ur{لنگوٹ}}
\textbf{\ur{لنگوٹ}} — \emph{langoṭ} — \textquotedblleft{}a loincloth (a wrestler's)\textquotedblright{}.

\begin{quote}
\ur{پہلوان} \ur{لنگوٹ} \ur{پہن} \ur{کر} \ur{میدان} \ur{میں} \ur{اترے۔}
\end{quote}

\begin{quote}
\emph{pahalvān langoṭ pahan kar maidān meṅ utre.}
\end{quote}

\begin{quote}
The wrestlers put on their loincloths and stepped into the ring.
\end{quote}

Masculine.

\begin{cousinweb}
\textbf{\ur{لنگوٹ}} (\emph{langoṭ}) is inherited: CDIAL lists the Hindi cousin \emph{lãgoṭ}, \textquotedblleft{}a loincloth\textquotedblright{}, under \emph{*laṅgapaṭṭa}, a form it reconstructs (CDIAL 10903).
\end{cousinweb}

\begin{grammarlens}[title={What you've built}]
One more word for describing what you lived and felt.
\end{grammarlens}

\subsection*{Your turn}
\begin{itemize}
\raggedright
  \item \emph{Say it:} \emph{langoṭ}
  \item \emph{Say it:} \emph{langoṭ}, once more
  \item \emph{Say it:} \emph{pahalvān langoṭ pahan kar maidān meṅ utre.} — The wrestlers put on their loincloths and stepped into the ring.
  \item \emph{Recall:} say the Urdu for a move (in wrestling), then the Urdu for to pit against each other, then say \emph{langoṭ} again
\end{itemize}

\begin{tcolorbox}[breakable,colback=teal!4,colframe=teal!35!black,title={Before you move on}]
What does \ur{لنگوٹ} mean? (\textquotedblleft{}A loincloth (a wrestler's)\textquotedblright{}.) Say it once more.
\end{tcolorbox}
